% Bagian: incompleteness
% Bab: theories-computability
% Subbagian: computably-axiomatizable

\documentclass[../../../include/open-logic-section]{subfiles}

\begin{document}

\olfileid[id]{inc}{tcp}{cax}

\olsection{Teori yang \printtoken{S}{axiomatizable}}

Suatu teori~$\Th{T}$ disebut \emph{!!{axiomatizable}} jika teori itu memiliki
suatu himpunan aksioma~$A$ yang dapat dihitung. (Pernyataan bahwa $A$ adalah
himpunan aksioma bagi~$\Th{T}$ berarti $T = \Setabs{!A}{A \Proves !A}$.)
Setiap aksiomatisasi bilangan asli yang ``wajar'' akan memiliki sifat ini.
Secara khusus, setiap teori dengan himpunan aksioma berhingga bersifat
!!{axiomatizable}.

\begin{lem}
Andaikan $\Th{T}$ bersifat !!{axiomatizable}. Maka $\Th{T}$
!!{computably enumerable}.
\end{lem}

\begin{proof}
Andaikan $A$ merupakan himpunan aksioma yang dapat dihitung bagi~$\Th{T}$.
Untuk menentukan apakah $!A \in T$, cukup cari !!a{derivation} atas~$!A$
dari aksioma-aksioma tersebut.

Dengan ungkapan yang sedikit berbeda, $!A$ berada dalam~$\Th{T}$ jika dan
hanya jika terdapat suatu daftar berhingga aksioma $!B_1$, \dots, $!B_k$
dalam~$A$ dan !!a{derivation} atas $(!B_1 \land \dots \land !B_k) \lif !A$
dalam logika orde pertama. Namun, kita telah mengetahui bahwa setiap himpunan
dengan definisi berbentuk ``terdapat \dots sedemikian sehingga \dots''
bersifat !!{c.e.}, asalkan bagian ``\dots'' yang kedua dapat dihitung.
\end{proof}

\end{document}
